\chapter{Другая четверть}\label{ch:chetverti}

\textbf{Что вы узнете:} как перенести готовое шоу в другую часть коврика.

\section{Сдвиг номеров}

Номера клеток всегда считаются по всему коврику, а четверть --- окно на его часть. Чтобы то же шоу полетело в другой четверти, нужно сдвинуть номера клеток. Левая нижняя клетка каждой четверти:

\begin{center}
\begin{tabular}{@{}llr@{}}
\toprule
Четверть & Где & Сдвиг номеров \\
\midrule
1 & слева внизу & 0 \\
2 & справа внизу & $+6$ \\
3 & слева вверху & $+72$ \\
4 & справа вверху & $+78$ \\
\bottomrule
\end{tabular}
\end{center}

\begin{figure}[H]
  \centering
  % Четыре четверти коврика и сдвиг номеров клеток
\begin{tikzpicture}[x=1.0cm, y=1.0cm, font=\sffamily]
  \foreach \q/\x/\y/\s in {1/0/0/0, 2/6/0/6, 3/0/6/72, 4/6/6/78}{
    \fill[accentsoft] (\x,\y) rectangle (\x+6,\y+6);
    \draw[accent, very thick] (\x,\y) rectangle (\x+6,\y+6);
    \node[text=accent, font=\large\bfseries] at (\x+3,\y+4.4) {Четверть \q};
    \node[text=ink, font=\small] at (\x+3,\y+3.4) {левая нижняя клетка: \s};
    \node[text=inkgray, font=\small] at (\x+3,\y+2.4) {сдвиг номеров: \ifnum\s=0 нет\else $+\s$\fi};
  }
  \node[text=inkgray, font=\scriptsize] at (6,-0.4) {клетка 0 внизу слева, номера растут вправо и вверх};
\end{tikzpicture}

  \caption{Четыре четверти коврика и сдвиг номеров клеток.}\label{fig:quarters}
\end{figure}

Вот квадрат из урока~\ref{ch:marshrut}, перенесённый во вторую четверть. Ко всем клеткам прибавлено 6: 13 $\to$ 19, 16 $\to$ 22, 52 $\to$ 58, 49 $\to$ 55.

\solfile{l14_quarter2.py}

\section{Перенос в меню «Шоу…»}

Складывать вручную не обязательно. В меню \knopka{Шоу…} над сценой есть раздел «Перенести шоу в четверть»: выберите четверть, подтвердите «Сдвинуть шоу в четверть N? Номера клеток в коде изменятся», и номера клеток сдвинутся сами. Отменить можно клавишами \klavisha{Ctrl}+\klavisha{Z} в редакторе.

\begin{vnimanie}[Область и перенос]
Выбор области сам код не меняет, это только рамка для проверки. Если шоу написано для клеток 0--5, а область выбрана другая, проверка скажет «полёт выходит за рабочую область». Тогда смените область или перенесите шоу.
\end{vnimanie}

\begin{probuy}
  \item Перенесите квадрат из урока~\ref{ch:marshrut} в третью четверть (прибавьте 72).
  \item Перенесите тот же квадрат в четвёртую четверть (прибавьте 78).
  \item Чему равна клетка 14 в третьей четверти? А клетка 55 в четвёртой?
  \item Клетка 100: в какой она четверти? Какой номер у клетки на том же месте в четверти 1?
\end{probuy}
\otvety{14}
